\documentclass[11pt]{article}
\usepackage{mathblatt}


\begin{document}
\pruefheftstil{Quadratische}{Prüfungsheft P10 \textperiodcentered{} schwach}
\noindent{\Large\bfseries Quadratische}\hfill{\small\color{mbgrau}Prüfungsheft P10 \textperiodcentered{} schwach}\par\medskip
\pfrueckblick
\begin{pfnr}{}{1.}{}
Das Koordinatensystem zeigt den Punkt $A$. Gib seine Koordinaten an.
\par\smallskip\noindent \begin{ksys}[xmin=-4,xmax=4,ymin=-4,ymax=4,ablesen]
\punkt{3}{2}{A}
\end{ksys}\par
\par\noindent A(\leerfeld[|\_\_)] \par
\fusshilfe{\mbox{1:~$A(3|2)$}}
\end{pfnr}
\begin{pfnr}{}{2.}{}
Trage den Punkt $B(4|1)$ in das Koordinatensystem ein.
\par\smallskip\noindent \begin{ksys}[xmin=-4,xmax=5,ymin=-4,ymax=4]
\end{ksys}\par
\rechenraster{1}
\fusshilfe{\mbox{2:~$B$: $4$ nach rechts, $1$ nach oben}}
\end{pfnr}
\begin{pfnr}{}{3.}{}
Berechne $g(4)$ für $g(x) = 2x + 1$.
\par\noindent \leerfeld \par
\rechenraster{1}
\fusshilfe{\mbox{3:~$9$}}
\end{pfnr}
\begin{pfnr}{}{4.}{}
Zeichne die Gerade $g(x) = x + 2$ in das Koordinatensystem.
\par\smallskip\noindent \begin{ksys}[xmin=-4,xmax=4,ymin=-3,ymax=6]
\end{ksys}\par
\rechenraster{1}
\fusshilfe{\mbox{4:~Gerade durch $(0|2)$ und $(2|4)$}}
\end{pfnr}
\begin{pfnr}{}{5.}{}
Multipliziere aus. $(x + 4)^2$
\par\noindent \leerfeld \par
\fusshilfe{\mbox{5:~$x^2 + 8x + 16$}}
\end{pfnr}
\begin{pfnr}{}{6.}{}
Multipliziere aus. $(x + 9)^2$
\par\noindent \leerfeld \par
\fusshilfe{\mbox{6:~$x^2 + 18x + 81$}}
\end{pfnr}
\begin{pfnr}{}{7.}{}
Löse die Gleichung. $3x + 2 = 14$
\par\noindent x = \leerfeld \par
\rechenraster{1}
\fusshilfe{\mbox{7:~$4$}}
\end{pfnr}
\begin{pfnr}{}{8.}{}
Löse die Gleichung. $x - 7 = 5$
\par\noindent x = \leerfeld \par
\rechenraster{1}
\fusshilfe{\mbox{8:~$12$}}
\end{pfnr}
\begin{pfnr}{}{9.}{}
Gegeben sind die Geraden $g(x) = x + 1$ und $h(x) = -x + 5$. Berechne ihren Schnittpunkt.
\par\noindent S(\leerfeld[|\_\_)] \par
\rechenraster{1}
\fusshilfe{\mbox{9:~$S(2|3)$}}
\end{pfnr}
\begin{pfnr}{}{10.}{}
Gegeben sind die Geraden $g(x) = 2x$ und $h(x) = x + 3$. Berechne ihren Schnittpunkt.
\par\noindent S(\leerfeld[|\_\_)] \par
\rechenraster{1}
\fusshilfe{\mbox{10:~$S(3|6)$}}
\end{pfnr}
\begin{pfnr}{}{11.}{}
Berechne. $\sqrt{64}$
\par\noindent \leerfeld \par
\fusshilfe{\mbox{11:~$8$}}
\end{pfnr}
\begin{pfnr}{}{12.}{}
Berechne. $\sqrt{121}$
\par\noindent \leerfeld \par
\fusshilfe{\mbox{12:~$11$}}
\end{pfnr}
\begin{pfnr}{}{13.}{}
Löse die Gleichung $x^2 = 36$. Gib die Lösungsmenge an.
\par\noindent x\_1 = \leerfeld , x\_2 = \leerfeld \par
\fusshilfe{\mbox{13:~$x_1 = 6$, $x_2 = -6$}}
\end{pfnr}
\begin{pfnr}{}{14.}{}
Löse die Gleichung $x^2 = 100$. Gib die Lösungsmenge an.
\par\noindent x\_1 = \leerfeld , x\_2 = \leerfeld \par
\fusshilfe{\mbox{14:~$x_1 = 10$, $x_2 = -10$}}
\end{pfnr}
\begin{pfnr}{}{15.}{}
Berechne. $3 \cdot (2 + 4)$
\par\noindent \leerfeld \par
\rechenraster{1}
\fusshilfe{\mbox{15:~$18$}}
\end{pfnr}
\begin{pfnr}{}{16.}{}
Berechne. $(10 - 4) \cdot 2$
\par\noindent \leerfeld \par
\rechenraster{1}
\fusshilfe{\mbox{16:~$12$}}
\end{pfnr}
\begin{pfnr}{}{17.}{}
Multipliziere aus. $x \cdot (x + 4)$
\par\noindent \leerfeld \par
\fusshilfe{\mbox{17:~$x^2 + 4x$}}
\end{pfnr}
\begin{pfnr}{}{18.}{}
Klammere $x$ aus. $x^2 + 7x$
\par\noindent \leerfeld \par
\fusshilfe{\mbox{18:~$x \cdot (x + 7)$}}
\end{pfnr}
\pfstufekopf[sscheitelpunktablesen]{Scheitelpunkt ablesen}{in 2 der letzten 5 Prüfungen}
\pfgruppe{}
\begin{pfnr}{eigene Aufgabe}{19.}{}
Kreise in $f(x) = (x - 4)^2 + 2$ die Zahl in der Klammer ein. Bei welchem $x$ liegt der Scheitel?
\par\noindent x = \leerfeld \par
\rechenraster{1}
\fusshilfe{\mbox{19:~$4$}}
\end{pfnr}
\begin{pfnr}{eigene Aufgabe}{20.}{}
Kreuze an, in welcher Form $f(x) = x^2 + 4x - 1$ geschrieben ist.
\pfkreuzzeile{Scheitelpunktform}\pfkreuzzeile{Normalform}
\fusshilfe{\mbox{20:~Normalform}}
\end{pfnr}
\pfgruppe{}
\begin{pfnr}{eigene Aufgabe}{21.}{}
Gib den Scheitel der Parabel $f(x) = (x - 2)^2 + 5$ an.
\par\noindent S(\leerfeld[|\_\_)] \par
\fusshilfe{\mbox{21:~$S(2|5)$}}
\end{pfnr}
\begin{pfnr}{eigene Aufgabe}{22.}{}
Gib den Schnittpunkt der Parabel $f(x) = x^2 + 3x + 5$ mit der $y$-Achse an.
\par\noindent (\leerfeld[|\_\_)] \par
\rechenraster{1}
\fusshilfe{\mbox{22:~Schnittpunkt $(0|5)$}}
\end{pfnr}
\begin{pfnr}{eigene Aufgabe}{23.}{}
Gib den Scheitel der Parabel $f(x) = (x + 4)^2 + 2$ an.
\par\noindent S(\leerfeld[|\_\_)] \par
\fusshilfe{\mbox{23:~$S(-4|2)$}}
\end{pfnr}
\begin{pfnr}{eigene Aufgabe}{24.}{}
Das Koordinatensystem zeigt die Parabel $f(x) = x^2 - 4x + 1$. Lies ihren Scheitel ab.
\par\smallskip\noindent \begin{ksys}[xmin=-1,xmax=5,ymin=-4,ymax=2,ablesen]
\parabel{1}{2}{-3}{f}
\end{ksys}\par
\par\noindent S(\leerfeld[|\_\_)] \par
\fusshilfe{\mbox{24:~$S(2|{-3})$}}
\end{pfnr}
\begin{pfnr}{eigene Aufgabe}{25.}{}
Das Koordinatensystem zeigt eine Parabel. Lies ihren Scheitel ab.
\par\smallskip\noindent \begin{ksys}[xmin=-5,xmax=1,ymin=-1,ymax=8,ablesen]
\parabel{1}{-2}{3}{f}
\end{ksys}\par
\par\noindent S(\leerfeld[|\_\_)] \par
\fusshilfe{\mbox{25:~$S(-2|3)$}}
\end{pfnr}
\pfgruppe{}
\begin{pfnr}{P10 ’17}{26.}{1 BE}
Die Parabel $p$ hat die Gleichung $y = (x+3)^2 - 2$ und ist im Koordinatensystem gezeichnet. Gib den Scheitelpunkt der Parabel $p$ an: $S(\ \ |\ \ )$.
\par\smallskip\noindent \begin{tikzpicture}\begin{axis}[xmin=-6,xmax=2,ymin=-2,ymax=4,axis lines=middle,tick label style={font=\scriptsize},clip=true,/pgf/number format/use comma,axis line style={-{Stealth[length=2mm]}},every axis plot/.append style={line width=0.9pt},x=0.55cm,y=0.55cm,xtick distance=1,ytick distance=1,enlargelimits=false,grid=both,grid style={mbgitter,line width=0.3pt},major grid style={mbgitter},xlabel={$x$},ylabel={$y$},x label style={anchor=west},y label style={anchor=south}]\addplot[black,domain=-6:2,samples=120] {((x+3)^2)-2};\node[font=\small,fill=white,inner sep=1pt,anchor=south west] at (axis cs:-0.752,3.0535) {$p$};\end{axis}\end{tikzpicture}\par
\rechenraster{1}
\fusshilfe{\mbox{26:~S($-$3; $-$2)}}
\end{pfnr}
\begin{pfnr}{P10 ’18}{27.}{1 BE}
Im Koordinatensystem ist die Parabel $p$ mit $p(x) = x^2 - 4x + 2$ gezeichnet. Gib den Scheitelpunkt von $p$ an: $S(\ \ |\ \ )$.
\par\smallskip\noindent \begin{tikzpicture}\begin{axis}[xmin=-4,xmax=5,ymin=-2,ymax=7,axis lines=middle,tick label style={font=\scriptsize},clip=true,/pgf/number format/use comma,axis line style={-{Stealth[length=2mm]}},every axis plot/.append style={line width=0.9pt},x=0.55cm,y=0.55cm,xtick distance=1,ytick distance=1,enlargelimits=false,grid=both,grid style={mbgitter,line width=0.3pt},major grid style={mbgitter},xlabel={$x$},ylabel={$y$},x label style={anchor=west},y label style={anchor=south}]\addplot[black,domain=-4:5,samples=120] {(x^2)-4*x+2};\node[font=\small,fill=white,inner sep=1pt,anchor=south west] at (axis cs:4.28,3.1984) {$p$};\end{axis}\end{tikzpicture}\par
\rechenraster{1}
\fusshilfe{\mbox{27:~S(2; $-$2)}}
\end{pfnr}
\begin{pfnr}{P10 ’20}{28.}{1 BE}
Die Parabel zu $f$ mit $y = (x+2)^2 - 4$ ist im Koordinatensystem gezeichnet. $S$ ist der Scheitelpunkt des Graphen von $f$. Gib seine Koordinaten an.
\par\smallskip\noindent \begin{tikzpicture}\begin{axis}[xmin=-5,xmax=7,ymin=-5,ymax=5,axis lines=middle,tick label style={font=\scriptsize},clip=true,/pgf/number format/use comma,axis line style={-{Stealth[length=2mm]}},every axis plot/.append style={line width=0.9pt},x=0.55cm,y=0.55cm,xtick distance=1,ytick distance=1,enlargelimits=false,grid=both,grid style={mbgitter,line width=0.3pt},major grid style={mbgitter},xlabel={$x$},ylabel={$y$},x label style={anchor=west},y label style={anchor=south}]\addplot[black,domain=-5:7,samples=120] {((x+2)^2)-4};\node[font=\small,fill=white,inner sep=1pt,anchor=south west] at (axis cs:0.76,3.6176) {$f$};\end{axis}\end{tikzpicture}\par
\rechenraster{1}
\fusshilfe{\mbox{28:~S($-$2; $-$4)}}
\end{pfnr}
\begin{pfnr}{P10 ’22}{29.}{1 BE}
Die lineare Funktion $f$ hat die Gleichung $f(x) = -2x + 3$. Die verschobene Normalparabel $p$ hat die Gleichung $p(x) = (x-2)^2 - 4$. Gib den Scheitelpunkt $S$ von $p$ an.
\rechenraster{1}
\fusshilfe{\mbox{29:~S(2|$-$4)}}
\end{pfnr}
\pfgruppe{}
\begin{pfnr}{P10 ’21}{30.}{1 BE}
Kreuze den Scheitelpunkt der gezeichneten Parabel an:
\pfkreuzzeile{$S(-1\,|\,-2)$}\pfkreuzzeile{$S(-1\,|\,2)$}\pfkreuzzeile{$S(-2\,|\,-1)$}\pfkreuzzeile{$S(2\,|\,-1)$}
\par\smallskip\noindent \begin{tikzpicture}\begin{axis}[xmin=-4,xmax=2,ymin=-2,ymax=3,axis lines=middle,tick label style={font=\scriptsize},clip=true,/pgf/number format/use comma,axis line style={-{Stealth[length=2mm]}},every axis plot/.append style={line width=0.9pt},x=0.55cm,y=0.55cm,xtick distance=1,ytick distance=1,enlargelimits=false,grid=both,grid style={mbgitter,line width=0.3pt},major grid style={mbgitter},xlabel={$x$},ylabel={$y$},x label style={anchor=west},y label style={anchor=south}]\addplot[black,domain=-4:2,samples=120] {((x+2)^2)-1};\end{axis}\end{tikzpicture}\par
\fusshilfe{\mbox{30:~S($-$2; $-$1)}}
\end{pfnr}
\pfstufekopf[sscheitelpunktformang]{Scheitelpunktform angeben}{in 2 der letzten 5 Prüfungen}
\pfgruppe{}
\begin{pfnr}{eigene Aufgabe}{31.}{}
Eine Normalparabel ist nach oben geöffnet. Ihr Scheitel ist $S(5|2)$. Stelle ihre Gleichung auf.
\par\noindent f(x) = \leerfeld \par
\rechenraster{1}
\fusshilfe{\mbox{31:~$(x - 5)^2 + 2$}}
\end{pfnr}
\begin{pfnr}{eigene Aufgabe}{32.}{}
Bringe $f(x) = x^2 - 8x + 19$ mit quadratischer Ergänzung in die Scheitelpunktform. Gib den Scheitel an.
\par\noindent f(x) = \leerfeld , S(\leerfeld[|\_\_)] \par
\rechenraster{1}
\fusshilfe{\mbox{32:~$S(4|3)$}}
\end{pfnr}
\begin{pfnr}{eigene Aufgabe}{33.}{}
Das Koordinatensystem zeigt eine Parabel. Stelle ihre Gleichung auf.
\par\smallskip\noindent \begin{ksys}[xmin=-1,xmax=5,ymin=-4,ymax=2,ablesen]
\parabel{1}{2}{-3}{f}
\end{ksys}\par
\par\noindent f(x) = \leerfeld \par
\rechenraster{1}
\fusshilfe{\mbox{33:~$(x - 2)^2 - 3$}}
\end{pfnr}
\pfgruppe{}
\begin{pfnr}{eigene Aufgabe}{34.}{}
Die Parabel $f(x) = (x - 4)^2 + 2$ wird an der $y$-Achse gespiegelt. So entsteht das Spiegelbild $g$. Gib die Gleichung von $g$ an.
\par\noindent g(x) = \leerfeld \par
\rechenraster{1}
\fusshilfe{\mbox{34:~$(x + 4)^2 + 2$, Scheitel $S(-4|2)$}}
\end{pfnr}
\begin{pfnr}{eigene Aufgabe}{35.}{}
Die Parabel $f(x) = (x - 1)^2 + 2$ wird um $4$ nach unten verschoben. So entsteht die Parabel $g$. Gib die Gleichung von $g$ an.
\par\noindent g(x) = \leerfeld \par
\rechenraster{1}
\fusshilfe{\mbox{35:~$(x - 1)^2 - 2$, Scheitel $S(1|{-2})$}}
\end{pfnr}
\begin{pfnr}{eigene Aufgabe}{36.}{}
Die Parabel $f(x) = (x - 5)^2 + 1$ wird an der $x$-Achse gespiegelt. So entsteht das Spiegelbild $g$. Gib die Gleichung von $g$ an.
\par\noindent g(x) = \leerfeld \par
\rechenraster{1}
\fusshilfe{\mbox{36:~$-(x - 5)^2 - 1$, Scheitel $S(5|{-1})$}}
\end{pfnr}
\begin{pfnr}{eigene Aufgabe}{37.}{}
Das Koordinatensystem zeigt die Parabel $f(x) = x^2 - 6x + 5$. Lies ihren Scheitel ab. Stelle damit die Scheitelpunktform auf. Prüfe sie mit $x = 0$.
\par\smallskip\noindent \begin{ksys}[xmin=-1,xmax=6,ymin=-5,ymax=1,ablesen]
\parabel{1}{3}{-4}{f}
\end{ksys}\par
\par\noindent f(x) = \leerfeld \par
\rechenraster{2}
\fusshilfe{\mbox{37:~$(x - 3)^2 - 4$; $5$}}
\end{pfnr}
\begin{pfnr}{eigene Aufgabe}{38.}{}
Gesucht ist eine nach oben geöffnete Normalparabel. Ihr Scheitel soll auf der $x$-Achse rechts vom Ursprung liegen. Gib eine passende Gleichung an.
\par\noindent f(x) = \leerfeld \par
\rechenraster{1}
\fusshilfe{\mbox{38:~$(x - 4)^2$ mit $S(4|0)$}}
\end{pfnr}
\pfgruppe{}
\begin{pfnr}{P10 ’16}{39.}{1 BE}
Der Scheitelpunkt einer verschobenen Normalparabel ist $S(1\,|\,3)$. Kreuze die Gleichung an, die passt:
\pfkreuzzeile{$y = (x+1)^2 + 3$}\pfkreuzzeile{$y = (x-1)^2 - 3$}\pfkreuzzeile{$y = (x-1)^2 + 3$}
\fusshilfe{\mbox{39:~y = (x $-$ 1)² + 3}}
\end{pfnr}
\pfgruppe{}
\begin{pfnr}{P10 ’15}{\llap{\textasteriskcentered\,}40.}{1 BE}
Die gezeichnete Normalparabel wird um 2 Einheiten nach rechts verschoben. Gib den Scheitelpunkt der verschobenen Parabel an: $S(\ \ |\ \ )$.
\par\smallskip\noindent \begin{tikzpicture}\begin{axis}[xmin=-2,xmax=2,ymin=-1,ymax=4,axis lines=middle,tick label style={font=\scriptsize},clip=true,/pgf/number format/use comma,axis line style={-{Stealth[length=2mm]}},every axis plot/.append style={line width=0.9pt},x=0.55cm,y=0.55cm,xtick distance=1,ytick distance=1,enlargelimits=false,grid=both,grid style={mbgitter,line width=0.3pt},major grid style={mbgitter},xlabel={$x$},ylabel={$y$},x label style={anchor=west},y label style={anchor=south}]\addplot[black,domain=-2:2,samples=120] {(x^2)};\end{axis}\end{tikzpicture}\par
\rechenraster{1}
\fusshilfe{\mbox{40:~S(2; 0)}}
\end{pfnr}
\begin{pfnr}{P10 ’18}{\llap{\textasteriskcentered\,}41.}{2 BE}
Im Koordinatensystem ist die Parabel $p$ mit $p(x) = x^2 - 4x + 2$ gezeichnet. Schreib die Gleichung von $p$ in Scheitelpunktform.
\par\smallskip\noindent \begin{tikzpicture}\begin{axis}[xmin=-4,xmax=5,ymin=-2,ymax=7,axis lines=middle,tick label style={font=\scriptsize},clip=true,/pgf/number format/use comma,axis line style={-{Stealth[length=2mm]}},every axis plot/.append style={line width=0.9pt},x=0.55cm,y=0.55cm,xtick distance=1,ytick distance=1,enlargelimits=false,grid=both,grid style={mbgitter,line width=0.3pt},major grid style={mbgitter},xlabel={$x$},ylabel={$y$},x label style={anchor=west},y label style={anchor=south}]\addplot[black,domain=-4:5,samples=120] {(x^2)-4*x+2};\node[font=\small,fill=white,inner sep=1pt,anchor=south west] at (axis cs:4.28,3.1984) {$p$};\end{axis}\end{tikzpicture}\par
\rechenraster{1}
\fusshilfe{\mbox{41:~p(x) = (x $-$ 2)² $-$ 2}}
\end{pfnr}
\begin{pfnr}{P10 ’20}{\llap{\textasteriskcentered\,}42.}{3 BE}
Die Parabel zu $f$ mit $y = (x+2)^2 - 4$ ist im Koordinatensystem gezeichnet. Spiegelt man den Graphen von $f$ an der y-Achse, entsteht die Parabel $p$. Skizziere $p$ in das Koordinatensystem. Gib eine Gleichung von $p$ an.
\par\smallskip\noindent \begin{tikzpicture}\begin{axis}[xmin=-5,xmax=7,ymin=-5,ymax=5,axis lines=middle,tick label style={font=\scriptsize},clip=true,/pgf/number format/use comma,axis line style={-{Stealth[length=2mm]}},every axis plot/.append style={line width=0.9pt},x=0.55cm,y=0.55cm,xtick distance=1,ytick distance=1,enlargelimits=false,grid=both,grid style={mbgitter,line width=0.3pt},major grid style={mbgitter},xlabel={$x$},ylabel={$y$},x label style={anchor=west},y label style={anchor=south}]\addplot[black,domain=-5:7,samples=120] {((x+2)^2)-4};\node[font=\small,fill=white,inner sep=1pt,anchor=south west] at (axis cs:0.76,3.6176) {$f$};\end{axis}\end{tikzpicture}\par
\rechenraster{1}
\fusshilfe{\mbox{42:~Skizze: Normalparabel mit Scheitel (2; $-$4); y = (x $-$ 2)² $-$ 4}}
\end{pfnr}
\begin{pfnr}{P10 ’17}{\llap{\textasteriskcentered\,}43.}{2 BE}
Die Parabel $p$ hat die Gleichung $y = (x+3)^2 - 2$ und ist im Koordinatensystem gezeichnet. Man verschiebt $p$ um 2 Einheiten nach oben und spiegelt das Ergebnis danach an der x-Achse. So entsteht die Parabel $q$. Gib eine Gleichung von $q$ an.
\par\smallskip\noindent \begin{tikzpicture}\begin{axis}[xmin=-6,xmax=2,ymin=-2,ymax=4,axis lines=middle,tick label style={font=\scriptsize},clip=true,/pgf/number format/use comma,axis line style={-{Stealth[length=2mm]}},every axis plot/.append style={line width=0.9pt},x=0.55cm,y=0.55cm,xtick distance=1,ytick distance=1,enlargelimits=false,grid=both,grid style={mbgitter,line width=0.3pt},major grid style={mbgitter},xlabel={$x$},ylabel={$y$},x label style={anchor=west},y label style={anchor=south}]\addplot[black,domain=-6:2,samples=120] {((x+3)^2)-2};\node[font=\small,fill=white,inner sep=1pt,anchor=south west] at (axis cs:-0.752,3.0535) {$p$};\end{axis}\end{tikzpicture}\par
\pfzwfrage{Gib die Gleichung der nach oben verschobenen Parabel an.}
\rechenraster{2}
\fusshilfe{\mbox{43:~q: y = $-$(x + 3)²}}
\end{pfnr}
\begin{pfnr}{P10 ’25}{44.}{2 BE}
Die quadratische Funktion $p$ hat die Gleichung $p(x) = x^2 - 6x + 7$. Ihr Graph ist im Koordinatensystem gezeichnet. Gib den Scheitelpunkt $S$ von $p$ an: $S(\ \ |\ \ )$. Schreib die Gleichung von $p$ in Scheitelpunktform: $p(x) = \ldots$
\par\smallskip\noindent \begin{tikzpicture}\begin{axis}[xmin=-3,xmax=6,ymin=-2,ymax=4,axis lines=middle,tick label style={font=\scriptsize},clip=true,/pgf/number format/use comma,axis line style={-{Stealth[length=2mm]}},every axis plot/.append style={line width=0.9pt},x=0.55cm,y=0.55cm,xtick distance=1,ytick distance=1,enlargelimits=false,grid=both,grid style={mbgitter,line width=0.3pt},major grid style={mbgitter},xlabel={$x$},ylabel={$y$},x label style={anchor=west},y label style={anchor=south}]\addplot[black,domain=-3:6,samples=120] {(x^2)-6*x+7};\node[font=\small,fill=white,inner sep=1pt,anchor=south west] at (axis cs:5.28,3.1984) {$p$};\end{axis}\end{tikzpicture}\par
\rechenraster{2}
\fusshilfe{\mbox{44:~S(3|$-$2); p(x) = (x $-$ 3)² $-$ 2}}
\end{pfnr}
\begin{pfnr}{P10 ’23}{\llap{\textasteriskcentered\,}45.}{2 BE}
Gib den Scheitelpunkt der gezeichneten Normalparabel $p$ an: $S(\ \ |\ \ )$. Schreib eine Gleichung von $p$ in Scheitelpunktform.
\par\smallskip\noindent \begin{tikzpicture}\begin{axis}[xmin=-5,xmax=5,ymin=-4,ymax=6,axis lines=middle,tick label style={font=\scriptsize},clip=true,/pgf/number format/use comma,axis line style={-{Stealth[length=2mm]}},every axis plot/.append style={line width=0.9pt},x=0.55cm,y=0.55cm,xtick distance=1,ytick distance=1,enlargelimits=false,grid=both,grid style={mbgitter,line width=0.3pt},major grid style={mbgitter},xlabel={$x$},ylabel={$y$},x label style={anchor=west},y label style={anchor=south}]\addplot[black,domain=-5:5,samples=120] {-((x+1)^2)+6};\node[font=\small,fill=white,inner sep=1pt,anchor=south west] at (axis cs:2,-3) {$p$};\end{axis}\end{tikzpicture}\par
\rechenraster{1}
\fusshilfe{\mbox{45:~S($-$1|6); p(x) = $-$(x + 1)² + 6}}
\end{pfnr}
\begin{pfnr}{P10 ’18}{\llap{\textasteriskcentered\,}46.}{4 BE}
Eine Normalparabel $q$ hat ihren Scheitelpunkt auf der y-Achse und keine Nullstelle. Skizziere eine solche Parabel in das Koordinatensystem und gib ihre Gleichung an. Spiegelt man $q$ an der x-Achse, entsteht $q^{\prime}$. Gib eine Gleichung von $q^{\prime}$ an.
\par\smallskip\noindent \begin{tikzpicture}\begin{axis}[xmin=-4,xmax=5,ymin=-6,ymax=6,axis lines=middle,tick label style={font=\scriptsize},clip=true,/pgf/number format/use comma,axis line style={-{Stealth[length=2mm]}},every axis plot/.append style={line width=0.9pt},x=0.55cm,y=0.55cm,xtick distance=1,ytick distance=1,enlargelimits=false,grid=both,grid style={mbgitter,line width=0.3pt},major grid style={mbgitter},xlabel={$x$},ylabel={$y$},x label style={anchor=west},y label style={anchor=south}]\end{axis}\end{tikzpicture}\par
\pfzwfrage{Gib an, ob der Scheitel über oder unter der x-Achse liegen muss.}
\rechenraster{2}
\fusshilfe{\mbox{46:~$-$x² $-$ 1}}
\end{pfnr}
\pfstufekopf[snullstellenberechnen]{Nullstellen berechnen}{in 1 der letzten 5 Prüfungen}
\pfgruppe{}
\begin{pfnr}{eigene Aufgabe}{47.}{}
Kreuze an, welche Form die Gleichung $x^2 - 3x + 1 = 0$ hat und welcher Lösungsweg passt.
\pfkreuzzeile{reinquadratisch – Wurzelziehen}\pfkreuzzeile{Produkt gleich null – Nullprodukt}\pfkreuzzeile{quadrierte Klammer – Rückwärtsrechnen}\pfkreuzzeile{Normalform – Formel}
\fusshilfe{\mbox{47:~Formel}}
\end{pfnr}
\begin{pfnr}{eigene Aufgabe}{48.}{}
Die Gleichung $x^2 - x - 12 = 0$ hat die Form $x^2 + px + q = 0$. Gib $p$ und $q$ an.
\par\noindent p = \leerfeld , q = \leerfeld \par
\rechenraster{1}
\fusshilfe{\mbox{48:~$-12$}}
\end{pfnr}
\begin{pfnr}{eigene Aufgabe}{49.}{}
Du willst die Schnittpunkte der Parabel $f(x) = x^2 - 2x$ und der Geraden $g(x) = x + 4$ berechnen. Schreibe nur die Gleichung auf, mit der du anfängst.
\par\noindent \leerfeld[=]  \leerfeld \par
\rechenraster{1}
\fusshilfe{\mbox{49:~$x + 4$}}
\end{pfnr}
\pfgruppe{}
\begin{pfnr}{eigene Aufgabe}{50.}{}
Berechne die Nullstellen von $f(x) = (x - 1)^2 - 4$.
\par\noindent x\_1 = \leerfeld , x\_2 = \leerfeld \par
\rechenraster{1}
\fusshilfe{\mbox{50:~$x_1 = 3$, $x_2 = -1$}}
\end{pfnr}
\begin{pfnr}{eigene Aufgabe}{51.}{}
Löse die Gleichung $x^2 + 6x + 8 = 0$. Gib die Lösungsmenge an.
\par\noindent x1 = \leerfeld , x2 = \leerfeld \par
\rechenraster{1}
\fusshilfe{\mbox{51:~$x_1 = -2$, $x_2 = -4$}}
\end{pfnr}
\begin{pfnr}{eigene Aufgabe}{52.}{}
Berechne die Nullstellen von $f(x) = x^2 - 2x - 8$.
\par\noindent x\_1 = \leerfeld , x\_2 = \leerfeld \par
\rechenraster{2}
\fusshilfe{\mbox{52:~$x_1 = 4$, $x_2 = -2$}}
\end{pfnr}
\begin{pfnr}{eigene Aufgabe}{53.}{}
Berechne die Nullstellen von $f(x) = 3x^2 - 3x - 18$.
\par\noindent x\_1 = \leerfeld , x\_2 = \leerfeld \par
\rechenraster{2}
\fusshilfe{\mbox{53:~$x_1 = 3$, $x_2 = -2$}}
\end{pfnr}
\begin{pfnr}{eigene Aufgabe}{54.}{}
Berechne die Nullstellen von $f(x) = x^2 - 4x + 1$. Runde auf zwei Stellen nach dem Komma.
\par\noindent x\_1 = \leerfeld , x\_2 = \leerfeld \par
\rechenraster{1}
\fusshilfe{\mbox{54:~$\approx 0{,}27$}}
\end{pfnr}
\pfgruppe{}
\begin{pfnr}{eigene Aufgabe}{55.}{}
Berechne die Nullstellen von $f(x) = x^2 + 12x + 36$. Wo liegt der Scheitel der Parabel?
\par\noindent x = \leerfeld , S(\leerfeld[|\_\_)] \par
\rechenraster{2}
\fusshilfe{\mbox{55:~$-6 \pm 0$; $-6$}}
\end{pfnr}
\begin{pfnr}{eigene Aufgabe}{56.}{}
Die Parabel $f(x) = x^2 - 4x + 7$ hat den Scheitel $S(2|3)$. Berechne ihre Nullstellen mit der p-q-Formel. Wie passt das Ergebnis zum Scheitel?
\rechenraster{3}
\end{pfnr}
\pfgruppe{}
\begin{pfnr}{P10 ’20}{\llap{\textasteriskcentered\,}57.}{3 BE}
Berechne die Nullstellen der Funktion mit $y = 2x^2 + 8x + 6$.
\pfzwfrage{Teile die Gleichung $2x^2 + 8x + 6 = 0$ durch 2.}
\rechenraster{3}
\fusshilfe{\mbox{57:~x$_1$ = $-$1; x$_2$ = $-$3}}
\end{pfnr}
\pfgruppe{}
\begin{pfnr}{P10 ’25}{\llap{\textasteriskcentered\,}58.}{4 BE}
Die quadratische Funktion $p$ hat die Gleichung $p(x) = x^2 - 6x + 7$. Ihr Graph ist im Koordinatensystem gezeichnet. Berechne die Koordinaten der Punkte, in denen die Parabel $p$ die x-Achse schneidet.
\par\smallskip\noindent \begin{tikzpicture}\begin{axis}[xmin=-3,xmax=6,ymin=-2,ymax=4,axis lines=middle,tick label style={font=\scriptsize},clip=true,/pgf/number format/use comma,axis line style={-{Stealth[length=2mm]}},every axis plot/.append style={line width=0.9pt},x=0.55cm,y=0.55cm,xtick distance=1,ytick distance=1,enlargelimits=false,grid=both,grid style={mbgitter,line width=0.3pt},major grid style={mbgitter},xlabel={$x$},ylabel={$y$},x label style={anchor=west},y label style={anchor=south}]\addplot[black,domain=-3:6,samples=120] {(x^2)-6*x+7};\node[font=\small,fill=white,inner sep=1pt,anchor=south west] at (axis cs:5.28,3.1984) {$p$};\end{axis}\end{tikzpicture}\par
\pfzwfrage{Setze $x^2 - 6x + 7 = 0$.}
\rechenraster{3}
\fusshilfe{\mbox{58:~N1(1,59|0), N2(4,41|0), x = 3 ± √2}}
\end{pfnr}
\begin{pfnr}{P10 ’17}{\llap{\textasteriskcentered\,}59.}{4 BE}
Die Parabel $p$ hat die Gleichung $y = (x+3)^2 - 2$ und ist im Koordinatensystem gezeichnet. Zeige, dass man die Gleichung von $p$ auch als $y = x^2 + 6x + 7$ schreiben kann. Berechne die Nullstellen von $p$.
\par\smallskip\noindent \begin{tikzpicture}\begin{axis}[xmin=-6,xmax=2,ymin=-2,ymax=4,axis lines=middle,tick label style={font=\scriptsize},clip=true,/pgf/number format/use comma,axis line style={-{Stealth[length=2mm]}},every axis plot/.append style={line width=0.9pt},x=0.55cm,y=0.55cm,xtick distance=1,ytick distance=1,enlargelimits=false,grid=both,grid style={mbgitter,line width=0.3pt},major grid style={mbgitter},xlabel={$x$},ylabel={$y$},x label style={anchor=west},y label style={anchor=south}]\addplot[black,domain=-6:2,samples=120] {((x+3)^2)-2};\node[font=\small,fill=white,inner sep=1pt,anchor=south west] at (axis cs:-0.752,3.0535) {$p$};\end{axis}\end{tikzpicture}\par
\rechenraster{4}
\fusshilfe{\mbox{59:~x² + 6x + 7; $-$3 $-$ √2 $\approx$ $-$4,41}}
\end{pfnr}
\pfstufekopf[sgeradeundparabelglei]{Gerade und Parabel gleichsetzen}{in 2 der letzten 5 Prüfungen}
\pfgruppe{}
\begin{pfnr}{eigene Aufgabe}{60.}{}
Gegeben sind die Parabel $f(x) = x^2 - 4x + 1$ und die Gerade $g(x) = 2x - 4$. Berechne ihre Schnittpunkte.
\par\noindent S\_1(\leerfeld[|\_\_)] , S\_2(\leerfeld[|\_\_)] \par
\rechenraster{2}
\fusshilfe{\mbox{60:~$S_1(5|6)$, $S_2(1|{-2})$}}
\end{pfnr}
\begin{pfnr}{eigene Aufgabe}{61.}{}
Berechne die gemeinsamen Punkte der Parabel $f(x) = x^2 + 1$ und der Geraden $g(x) = 2x$.
\par\noindent (\leerfeld[|\_\_)] \par
\rechenraster{3}
\fusshilfe{\mbox{61:~$2$}}
\end{pfnr}
\begin{pfnr}{eigene Aufgabe}{62.}{4 BE}
Berechne die gemeinsamen Punkte der Parabel $f(x) = x^2 + 4$ und der Geraden $g(x) = 2x$.
\rechenraster{3}
\fusshilfe{\mbox{62:~$-3$, negativ}}
\end{pfnr}
\begin{pfnr}{eigene Aufgabe}{63.}{}
Gegeben sind die Parabel $y = x^2$ und die Gerade $y = x + 6$. An welchen Stellen $x$ schneiden sie sich?
\par\noindent x1 = \leerfeld , x2 = \leerfeld \par
\rechenraster{3}
\fusshilfe{\mbox{63:~$x_1 = 3$, $x_2 = -2$}}
\end{pfnr}
\pfgruppe{}
\begin{pfnr}{eigene Aufgabe}{64.}{}
Gegeben sind die Parabel $f(x) = (x - 2)^2 - 1$ und die Gerade $g(x) = x - 3$. An welchen Stellen $x$ schneiden sie sich?
\par\noindent x\_1 = \leerfeld , x\_2 = \leerfeld \par
\rechenraster{2}
\fusshilfe{\mbox{64:~$x_1 = 3$, $x_2 = 2$}}
\end{pfnr}
\begin{pfnr}{eigene Aufgabe}{65.}{}
Gegeben sind die Parabel $f(x) = x^2 + x - 4$ und die Gerade $g(x) = 2x + 2$. An welchen Stellen $x$ schneiden sie sich?
\par\noindent x\_1 = \leerfeld , x\_2 = \leerfeld \par
\rechenraster{2}
\fusshilfe{\mbox{65:~$x_1 = 3$, $x_2 = -2$}}
\end{pfnr}
\pfgruppe{}
\begin{pfnr}{P10 ’21}{\llap{\textasteriskcentered\,}66.}{4 BE}
Die Parabel mit $y = x^2 + 2x - 1$ und die Gerade mit $y = 3x + 1$ schneiden sich. Berechne die Koordinaten der Schnittpunkte.
\pfzwfrage{Setze die beiden Terme gleich und bring alles auf eine Seite.}
\pfzwfrage{Löse die quadratische Gleichung.}
\pfzwfrage{Setze die x-Werte in die Geradengleichung ein.}
\rechenraster{4}
\fusshilfe{\mbox{66:~S$_1$($-$1; $-$2); S$_2$(2; 7)}}
\end{pfnr}
\begin{pfnr}{P10 ’22}{\llap{\textasteriskcentered\,}67.}{4 BE}
Die lineare Funktion $f$ hat die Gleichung $f(x) = -2x + 3$. Die verschobene Normalparabel $p$ hat die Gleichung $p(x) = (x-2)^2 - 4$. Bestimme die Koordinaten der Punkte, in denen die Gerade $f$ die Parabel $p$ schneidet.
\pfzwfrage{Setze $-2x + 3 = (x-2)^2 - 4$ und bring alles auf eine Seite.}
\rechenraster{4}
\fusshilfe{\mbox{67:~S$_1$($-$1|5), S$_2$(3|$-$3)}}
\end{pfnr}
\begin{pfnr}{P10 ’24}{\llap{\textasteriskcentered\,}68.}{4 BE}
Die Gerade $f$ ist der Graph von $f(x) = 4x + 1$. Die Parabel $p$ ist der Graph von $p(x) = x^2 - 4$. Die Graphen von $f$ und $p$ haben zwei Schnittpunkte. Berechne ihre Koordinaten.
\rechenraster{4}
\fusshilfe{\mbox{68:~S1($-$1|$-$3), S2(5|21)}}
\end{pfnr}
\pfstufekopf[swertetabellezuordnen]{Wertetabelle zuordnen}{in 1 der letzten 5 Prüfungen}
\pfgruppe{}
\begin{pfnr}{eigene Aufgabe}{69.}{}
Kreuze an, ob der Graph von $f(x) = 4x - 1$ eine Gerade oder eine Parabel ist.
\pfkreuzzeile{Gerade}\pfkreuzzeile{Parabel}
\fusshilfe{\mbox{69:~Gerade}}
\end{pfnr}
\pfgruppe{}
\begin{pfnr}{eigene Aufgabe}{70.}{}
Fülle die Wertetabelle zu $f(x) = x^2$ aus.
\wertetabelle{x}{f(x)}{-2,-1,0,1,2}
\rechenraster{1}
\fusshilfe{\mbox{70:~$4$}}
\end{pfnr}
\begin{pfnr}{eigene Aufgabe}{71.}{}
Die Wertetabelle gehört zu $f(x) = a \cdot x^2$. \\ \wertetabelle[7,28,63]{x}{f(x)}{1,2,3} Bestimme die Zahl $a$.
\par\noindent a = \leerfeld \par
\rechenraster{1}
\fusshilfe{\mbox{71:~$7x^2$}}
\end{pfnr}
\pfgruppe{}
\begin{pfnr}{eigene Aufgabe}{72.}{}
Das Koordinatensystem zeigt eine Parabel. Kreuze an, welche Gleichung zu ihr gehört. 
\pfkreuzzeile{$f(x) = x^2 + 3$}\pfkreuzzeile{$f(x) = -x^2 - 3$}\pfkreuzzeile{$f(x) = x^2 - 3$}\pfkreuzzeile{$f(x) = -x^2 + 3$}
\par\smallskip\noindent \begin{ksys}[xmin=-4,xmax=4,ymin=-5,ymax=5,ablesen]
\funktion{-1*\x^2+(3)}{f}
\end{ksys}\par
\fusshilfe{\mbox{72:~$3$}}
\end{pfnr}
\begin{pfnr}{eigene Aufgabe}{73.}{}
Kreuze an, welche Tabelle zu $f(x) = 0{,}5x^2$ gehört. \\ A: \wertetabelle[2,0{,}5,0,0{,}5,2]{x}{y}{-2,-1,0,1,2} B: \wertetabelle[1,0{,}25,0,0{,}25,1]{x}{y}{-2,-1,0,1,2} C: \wertetabelle[4,1,0,1,4]{x}{y}{-2,-1,0,1,2}
\pfkreuzzeile{A}\pfkreuzzeile{B}\pfkreuzzeile{C}
\fusshilfe{\mbox{73:~$2$}}
\end{pfnr}
\pfgruppe{}
\begin{pfnr}{P10 ’26}{74.}{1 BE}
Kreuze die Wertetabelle an, die zu $y = 3x^2$ gehört.
\par\smallskip\noindent \begin{tabular}[t]{@{}c@{}}\renewcommand{\arraystretch}{1.3}\begin{tabular}{|c|c|c|c|c|c|}\hline $x$ & $-$2 & $-$1 & 0 & 1 & 2 \\ \hline $y$ & 12 & 3 & 0 & 3 & 12 \\ \hline\end{tabular}\\[3pt]{\small Tabelle 1\enspace$\square$}\end{tabular}\hspace{0.5cm}\begin{tabular}[t]{@{}c@{}}\renewcommand{\arraystretch}{1.3}\begin{tabular}{|c|c|c|c|c|c|}\hline $x$ & $-$2 & $-$1 & 0 & 1 & 2 \\ \hline $y$ & $-$6 & $-$3 & 0 & 3 & 6 \\ \hline\end{tabular}\\[3pt]{\small Tabelle 2\enspace$\square$}\end{tabular}\hspace{0.5cm}\begin{tabular}[t]{@{}c@{}}\renewcommand{\arraystretch}{1.3}\begin{tabular}{|c|c|c|c|c|c|}\hline $x$ & $-$2 & $-$1 & 0 & 1 & 2 \\ \hline $y$ & 4 & 1 & 0 & 1 & 4 \\ \hline\end{tabular}\\[3pt]{\small Tabelle 3\enspace$\square$}\end{tabular}\par
\fusshilfe{\mbox{74:~erste Tabelle (y = 12, 3, 0, 3, 12)}}
\end{pfnr}
\begin{pfnr}{P10 ’15}{\llap{\textasteriskcentered\,}75.}{3 BE}
Lena verlässt ein Flugzeug mit dem Fallschirm. Die ersten 20 Sekunden stürzt sie frei, der Schirm ist noch zu; ihre Höhe sinkt dabei von 2400 m auf 1200 m. Der Graph zeigt ihre Höhe $h$ (in Metern) zur Zeit $t$ (in Sekunden). Welche Gleichung kann den Graphen in den ersten 20 Sekunden beschreiben? Kreuze an und begründe:
\pfkreuzzeile{$h(t) = 3t^2 + 2400$}\pfkreuzzeile{$h(t) = 3t^2 - 2400$}\pfkreuzzeile{$h(t) = -3t^2 + 2400$}\pfkreuzzeile{$h(t) = -3t^2 - 2400$}
\par\smallskip\noindent \begin{tikzpicture}\begin{axis}[xmin=0,xmax=280,ymin=0,ymax=2400,tick label style={font=\scriptsize},clip=true,/pgf/number format/use comma,axis line style={-{Stealth[length=2mm]}},every axis plot/.append style={line width=0.9pt},width=11.5cm,height=7.5cm,grid=both,grid style={mbgitter,line width=0.3pt},minor tick num=1,axis lines=left,xlabel={Zeit t in Sekunden},ylabel={Höhe h in Metern},label style={font=\small},scaled ticks=false,/pgf/number format/1000 sep={\,}]\addplot[black,domain=0:20,samples=120] {-3*(x^2)+2400};\node[font=\small,fill=white,inner sep=1pt,anchor=south west] at (axis cs:18.4,1384.32) {$G$};\addplot[black,domain=20:260,samples=120] {-5*x+1300};\end{axis}\end{tikzpicture}\par
\rechenraster{1}
\fusshilfe{\mbox{75:~$-$3t² + 2400}}
\end{pfnr}
\pfstufekopf[spunktaufderparabelpr]{Punkt auf der Parabel prüfen}{in 1 der letzten 5 Prüfungen}
\pfgruppe{}
\begin{pfnr}{eigene Aufgabe}{76.}{}
Berechne $f(-6)$ für $f(x) = x^2$.
\par\noindent f(-6) = \leerfeld \par
\rechenraster{1}
\fusshilfe{\mbox{76:~$36$}}
\end{pfnr}
\begin{pfnr}{eigene Aufgabe}{77.}{}
Berechne $f(-3)$ für $f(x) = x^2 + 4x - 2$.
\par\noindent f(-3) = \leerfeld \par
\rechenraster{1}
\fusshilfe{\mbox{77:~$-5$}}
\end{pfnr}
\pfgruppe{}
\begin{pfnr}{eigene Aufgabe}{78.}{4 BE}
Gegeben sind die Parabel $f(x) = x^2 - 5$ und die Gerade $g(x) = 2x + 3$. Prüfe, ob $P(4|11)$ ein Schnittpunkt der beiden Graphen ist.
\rechenraster{1}
\fusshilfe{\mbox{78:~11}}
\end{pfnr}
\pfgruppe{}
\begin{pfnr}{P10 ’20}{79.}{1 BE}
Die Parabel zu $f$ mit $y = (x+2)^2 - 4$ ist im Koordinatensystem gezeichnet. Genau einer dieser Punkte liegt nicht auf dem Graphen von $f$. Kreuze ihn an:
\pfkreuzzeile{$A(-3\,|\,-3)$}\pfkreuzzeile{$B(0\,|\,0)$}\pfkreuzzeile{$C(-1\,|\,-3)$}\pfkreuzzeile{$D(-1{,}5\,|\,-2{,}5)$}\pfkreuzzeile{$E(-2\,|\,-4)$}
\par\smallskip\noindent \begin{tikzpicture}\begin{axis}[xmin=-5,xmax=7,ymin=-5,ymax=5,axis lines=middle,tick label style={font=\scriptsize},clip=true,/pgf/number format/use comma,axis line style={-{Stealth[length=2mm]}},every axis plot/.append style={line width=0.9pt},x=0.55cm,y=0.55cm,xtick distance=1,ytick distance=1,enlargelimits=false,grid=both,grid style={mbgitter,line width=0.3pt},major grid style={mbgitter},xlabel={$x$},ylabel={$y$},x label style={anchor=west},y label style={anchor=south}]\addplot[black,domain=-5:7,samples=120] {((x+2)^2)-4};\node[font=\small,fill=white,inner sep=1pt,anchor=south west] at (axis cs:0.76,3.6176) {$f$};\end{axis}\end{tikzpicture}\par
\fusshilfe{\mbox{79:~D($-$1,5; $-$2,5)}}
\end{pfnr}
\begin{pfnr}{P10 ’18}{80.}{3 BE}
Im Koordinatensystem ist die Parabel $p$ mit $p(x) = x^2 - 4x + 2$ gezeichnet. Entscheide für jede Aussage, ob sie wahr oder falsch ist:
\par\smallskip\noindent \begingroup\renewcommand{\arraystretch}{1.5}\begin{tabular}{|p{8.0cm}|c|c|}\hline Aussage & wahr & falsch \\ \hline Der Punkt $(-1\,|\,7)$ liegt auf $p$. & $\square$ & $\square$ \\ \hline $p$ ist die Normalparabel, um 2 nach rechts und um 2 nach unten verschoben. & $\square$ & $\square$ \\ \hline $p$ schneidet die y-Achse im Punkt $(2\,|\,0)$ & $\square$ & $\square$ \\ \hline\end{tabular}\endgroup\par
\pfzwfrage{Berechne $p(-1)$.}
\pfzwfrage{Berechne $p(0)$.}
\fusshilfe{\mbox{80:~wahr; falsch}}
\end{pfnr}
\pfgruppe{}
\begin{pfnr}{P10 ’24}{81.}{2 BE}
Die Gerade $f$ ist der Graph von $f(x) = 4x + 1$. Die Parabel $p$ ist der Graph von $p(x) = x^2 - 4$. Prüfe mit einer Rechnung, ob der Punkt $T(-5\,|\,20)$ auf der Parabel $p$ liegt.
\rechenraster{1}
\fusshilfe{\mbox{81:~Nein, p($-$5) = 25 $-$ 4 = 21 $\neq$ 20}}
\end{pfnr}
\begin{pfnr}{P10 ’20}{82.}{2 BE}
Die Parabel zu $f$ mit $y = (x+2)^2 - 4$ ist im Koordinatensystem gezeichnet. Der Punkt $A(2\,|\,y)$ liegt auf dem Graphen von $f$ mit $y = (x+2)^2 - 4$. Berechne $y$.
\par\smallskip\noindent \begin{tikzpicture}\begin{axis}[xmin=-5,xmax=7,ymin=-5,ymax=5,axis lines=middle,tick label style={font=\scriptsize},clip=true,/pgf/number format/use comma,axis line style={-{Stealth[length=2mm]}},every axis plot/.append style={line width=0.9pt},x=0.55cm,y=0.55cm,xtick distance=1,ytick distance=1,enlargelimits=false,grid=both,grid style={mbgitter,line width=0.3pt},major grid style={mbgitter},xlabel={$x$},ylabel={$y$},x label style={anchor=west},y label style={anchor=south}]\addplot[black,domain=-5:7,samples=120] {((x+2)^2)-4};\node[font=\small,fill=white,inner sep=1pt,anchor=south west] at (axis cs:0.76,3.6176) {$f$};\end{axis}\end{tikzpicture}\par
\rechenraster{2}
\fusshilfe{\mbox{82:~y = 12}}
\end{pfnr}
\begin{pfnr}{P10 ’14}{\llap{\textasteriskcentered\,}83.}{4 BE}
Im Koordinatensystem sind die Parabel $p$ mit $p(x) = -x^2$ und die Gerade $g$ mit $g(x) = 2x - 3$ gezeichnet. Untersuche mit einer Rechnung, ob $S(-3\,|\,-9)$ ein gemeinsamer Punkt von $p$ und $g$ ist.
\par\smallskip\noindent \begin{tikzpicture}\begin{axis}[xmin=-6,xmax=6,ymin=-6,ymax=6,axis lines=middle,tick label style={font=\scriptsize},clip=true,/pgf/number format/use comma,axis line style={-{Stealth[length=2mm]}},every axis plot/.append style={line width=0.9pt},x=0.55cm,y=0.55cm,xtick distance=1,ytick distance=1,enlargelimits=false,grid=both,grid style={mbgitter,line width=0.3pt},major grid style={mbgitter},xlabel={$x$},ylabel={$y$},x label style={anchor=west},y label style={anchor=south}]\addplot[black,domain=-6:6,samples=120] {-(x^2)};\node[font=\small,fill=white,inner sep=1pt,anchor=south west] at (axis cs:2.136,-4.5625) {$p$};\addplot[black,domain=-6:6,samples=120] {2*x-3};\node[font=\small,fill=white,inner sep=1pt,anchor=south west] at (axis cs:3.72,4.44) {$g$};\end{axis}\end{tikzpicture}\par
\pfzwfrage{Berechne $p(-3)$.}
\rechenraster{2}
\fusshilfe{\mbox{83:~ja, S($-$3; $-$9) ist Schnittpunkt}}
\end{pfnr}
\pfstufekopf[sparabelskizzieren]{Parabel skizzieren}{in 1 der letzten 5 Prüfungen}
\pfgruppe{}
\begin{pfnr}{eigene Aufgabe}{84.}{}
Skizziere die Parabel zu $f(x) = (x - 2)^2 + 1$ im Koordinatensystem.
\par\smallskip\noindent \begin{ksys}[xmin=-2,xmax=6,ymin=-1,ymax=10]
\end{ksys}\par
\rechenraster{1}
\end{pfnr}
\pfgruppe{}
\begin{pfnr}{eigene Aufgabe}{85.}{}
Zeichne die Parabel zu $f(x) = 4x^2$ in das Koordinatensystem. Nimm dazu die Werte der Normalparabel und nimm sie mal vier.
\par\smallskip\noindent \begin{ksys}[xmin=-3,xmax=3,ymin=-2,ymax=18,ystep=2]
\end{ksys}\par
\rechenraster{1}
\end{pfnr}
\pfgruppe{}
\begin{pfnr}{P10 ’24}{86.}{2 BE}
Die Gerade $f$ ist der Graph von $f(x) = 4x + 1$. Die Parabel $p$ ist der Graph von $p(x) = x^2 - 4$. Skizziere die Parabel $p$ in das Koordinatensystem. Gib ihren Scheitelpunkt an: $S(\ \ |\ \ )$.
\par\smallskip\noindent \begin{tikzpicture}\begin{axis}[xmin=-4,xmax=4,ymin=-5,ymax=8,axis lines=middle,tick label style={font=\scriptsize},clip=true,/pgf/number format/use comma,axis line style={-{Stealth[length=2mm]}},every axis plot/.append style={line width=0.9pt},x=0.42cm,y=0.42cm,xtick distance=1,ytick distance=1,enlargelimits=false,grid=both,grid style={mbgitter,line width=0.3pt},major grid style={mbgitter},xlabel={$x$},ylabel={$y$},x label style={anchor=west},y label style={anchor=south}]\end{axis}\end{tikzpicture}\par
\rechenraster{1}
\fusshilfe{\mbox{86:~S(0|$-$4)}}
\end{pfnr}
\pfstufekopf[slagezweierparabelnoh]{Lage zweier Parabeln ohne Rechnung begründen}{in 1 der letzten 5 Prüfungen}
\pfgruppe{}
\begin{pfnr}{eigene Aufgabe}{87.}{}
Gegeben ist die Parabel $f(x) = x^2 - 2$. Gib eine Gerade $g$ an, die mit der Parabel keinen Punkt gemeinsam hat.
\par\noindent g(x) = \leerfeld \par
\rechenraster{1}
\end{pfnr}
\pfgruppe{}
\begin{pfnr}{eigene Aufgabe}{88.}{}
Gegeben sind die Parabeln $f(x) = (x - 1)^2 + 2$ und $g(x) = (x - 1)^2 - 4$. Beschreibe, wie sie zueinander liegen. Begründe ohne Rechnung, ob sie gemeinsame Punkte haben.
\rechenraster{1}
\end{pfnr}
\begin{pfnr}{eigene Aufgabe}{89.}{}
Wie viele Nullstellen hat $f(x) = (x - 2)^2 + 3$? Begründe ohne Rechnung.
\rechenraster{1}
\fusshilfe{\mbox{89:~Scheitel über der $x$-Achse, nach oben geöffnet}}
\end{pfnr}
\pfgruppe{}
\begin{pfnr}{P10 ’14}{\llap{\textasteriskcentered\,}90.}{2 BE}
Im Koordinatensystem sind die Parabel $p$ mit $p(x) = -x^2$ und die Gerade $g$ mit $g(x) = 2x - 3$ gezeichnet. Eine weitere Gerade $f$ soll die Parabel $p$ nirgends treffen. Gib eine Gleichung für eine solche Gerade an.
\par\smallskip\noindent \begin{tikzpicture}\begin{axis}[xmin=-6,xmax=6,ymin=-6,ymax=6,axis lines=middle,tick label style={font=\scriptsize},clip=true,/pgf/number format/use comma,axis line style={-{Stealth[length=2mm]}},every axis plot/.append style={line width=0.9pt},x=0.55cm,y=0.55cm,xtick distance=1,ytick distance=1,enlargelimits=false,grid=both,grid style={mbgitter,line width=0.3pt},major grid style={mbgitter},xlabel={$x$},ylabel={$y$},x label style={anchor=west},y label style={anchor=south}]\addplot[black,domain=-6:6,samples=120] {-(x^2)};\node[font=\small,fill=white,inner sep=1pt,anchor=south west] at (axis cs:2.136,-4.5625) {$p$};\addplot[black,domain=-6:6,samples=120] {2*x-3};\node[font=\small,fill=white,inner sep=1pt,anchor=south west] at (axis cs:3.72,4.44) {$g$};\end{axis}\end{tikzpicture}\par
\rechenraster{1}
\fusshilfe{\mbox{90:~z. B. f(x) = 1}}
\end{pfnr}
\pfstufekopf[slsungprfen]{Lösung prüfen}{in 1 der letzten 5 Prüfungen}
\pfgruppe{}
\begin{pfnr}{eigene Aufgabe}{91.}{}
Kreuze an, ob die Gleichung $5x - 2 = x^2$ quadratisch oder linear ist.
\pfkreuzzeile{quadratisch}\pfkreuzzeile{linear}
\fusshilfe{\mbox{91:~x kommt im Quadrat vor}}
\end{pfnr}
\begin{pfnr}{eigene Aufgabe}{92.}{}
Kreuze an, wie viele Lösungen die Gleichung $(x - 3)^2 = 25$ hat.
\pfkreuzzeile{zwei}\pfkreuzzeile{eine}\pfkreuzzeile{keine}
\fusshilfe{\mbox{92:~rechts positiv}}
\end{pfnr}
\begin{pfnr}{eigene Aufgabe}{93.}{}
Du willst $4x^2 + 8x - 12 = 0$ mit der Formel lösen. Kreuze an, ob und wodurch du vorher teilen musst.
\pfkreuzzeile{nicht teilen}\pfkreuzzeile{durch $4$}\pfkreuzzeile{durch $8$}
\fusshilfe{\mbox{93:~die Vorzahl von $x^2$}}
\end{pfnr}
\begin{pfnr}{eigene Aufgabe}{94.}{}
Der Satz vom Nullprodukt sagt: Ein Produkt ist null, wenn ein Faktor null ist. Kreuze an, ob der Satz für die Gleichung $(x - 2) \cdot (x + 7) = 0$ gilt.
\pfkreuzzeile{gilt}\pfkreuzzeile{gilt nicht}
\fusshilfe{\mbox{94:~rechts steht null}}
\end{pfnr}
\pfgruppe{}
\begin{pfnr}{eigene Aufgabe}{95.}{}
Löse die Gleichung $x^2 = 81$. Gib die Lösungsmenge an.
\par\noindent x1 = \leerfeld , x2 = \leerfeld \par
\fusshilfe{\mbox{95:~$x_1 = 9$, $x_2 = -9$}}
\end{pfnr}
\begin{pfnr}{eigene Aufgabe}{96.}{}
Löse die Gleichung $(x - 2) \cdot (x - 5) = 0$. Gib die Lösungsmenge an.
\par\noindent x1 = \leerfeld , x2 = \leerfeld \par
\rechenraster{1}
\fusshilfe{\mbox{96:~$x_1 = 2$, $x_2 = 5$}}
\end{pfnr}
\begin{pfnr}{eigene Aufgabe}{97.}{}
Prüfe, ob $x = -7$ eine Lösung von $(x + 2)^2 = 25$ ist.
\rechenraster{1}
\fusshilfe{\mbox{97:~$25$ : ja}}
\end{pfnr}
\begin{pfnr}{eigene Aufgabe}{98.}{}
Prüfe, ob $x = -5$ eine Lösung von $(x + 5) \cdot (x - 1) = 0$ ist.
\rechenraster{1}
\fusshilfe{\mbox{98:~$0$ : ja}}
\end{pfnr}
\begin{pfnr}{eigene Aufgabe}{99.}{3 BE}
Prüfe durch Einsetzen, ob $x = -2$ eine Lösung von $x^2 - 3x - 10 = 0$ ist.
\rechenraster{1}
\fusshilfe{\mbox{99:~$0$ : ja}}
\end{pfnr}
\pfgruppe{}
\begin{pfnr}{eigene Aufgabe}{100.}{}
Kreuze an, welcher Wert die Gleichung $x \cdot (x + 9) = -14$ erfüllt.
\pfkreuzzeile{$x = 2$}\pfkreuzzeile{$x = 7$}\pfkreuzzeile{$x = -7$}\pfkreuzzeile{$x = -14$}
\fusshilfe{\mbox{100:~$-14$}}
\end{pfnr}
\begin{pfnr}{P10 ’25}{101.}{1 BE}
Genau einer der Werte löst die Gleichung $x(x+5) = -6$. Kreuze ihn an:
\pfkreuzzeile{$x = 2$}\pfkreuzzeile{$x = 3$}\pfkreuzzeile{$x = -2$}\pfkreuzzeile{$x = -6$}
\fusshilfe{\mbox{101:~x = $-$2}}
\end{pfnr}
\pfstufekopf[sxzugegebenemy]{x zu gegebenem y}{in 1 der letzten 5 Prüfungen}
\pfgruppe{}
\begin{pfnr}{eigene Aufgabe}{102.}{}
Gegeben ist $f(x) = x^2 + x + 1$. Für welche $x$ ist $f(x) = 7$?
\par\noindent x\_1 = \leerfeld , x\_2 = \leerfeld \par
\rechenraster{2}
\fusshilfe{\mbox{102:~$x_1 = 2$, $x_2 = -3$}}
\end{pfnr}
\pfgruppe{}
\begin{pfnr}{eigene Aufgabe}{\llap{\textasteriskcentered\,}103.}{}
Die Parabel $f(x) = -x^2 + 4x + 2$ ist nach unten geöffnet. Für welche $x$ ist $f(x) = -10$?
\par\noindent x\_1 = \leerfeld , x\_2 = \leerfeld \par
\rechenraster{2}
\fusshilfe{\mbox{103:~$x_1 = 6$, $x_2 = -2$}}
\end{pfnr}
\begin{pfnr}{eigene Aufgabe}{\llap{\textasteriskcentered\,}104.}{}
Berechne die Nullstellen der Funktion $f$ mit $f(x) = 2x^2 - 2x - 40$.
\par\noindent x1 = \leerfeld , x2 = \leerfeld \par
\rechenraster{3}
\fusshilfe{\mbox{104:~$x_1 = 5$, $x_2 = -4$}}
\end{pfnr}
\pfgruppe{}
\begin{pfnr}{P10 ’23}{\llap{\textasteriskcentered\,}105.}{3 BE}
Die Parabel $p$ hat die Gleichung $p(x) = -x^2 - 2x + 5$. Bestimme die x-Koordinaten aller Punkte von $p$, deren y-Koordinate $-10$ ist.
\pfzwfrage{Setze $-x^2 - 2x + 5 = -10$ und bring alles auf eine Seite.}
\pfzwfrage{Multipliziere mit $-1$ und löse die quadratische Gleichung.}
\rechenraster{3}
\fusshilfe{\mbox{105:~x$_1$ = $-$5, x$_2$ = 3}}
\end{pfnr}
\pfpruefstein{P10 ’26}
\begin{pfnr}{}{}{}
Die verschobene Normalparabel $p$ ist der Graph von $p(x) = (x-2)^2 - 2$ und ist im Koordinatensystem gezeichnet.
\par\smallskip\noindent \begin{tikzpicture}\begin{axis}[xmin=-2,xmax=5,ymin=-2,ymax=5,axis lines=middle,tick label style={font=\scriptsize},clip=true,/pgf/number format/use comma,axis line style={-{Stealth[length=2mm]}},every axis plot/.append style={line width=0.9pt},x=0.55cm,y=0.55cm,xtick distance=1,ytick distance=1,enlargelimits=false,grid=both,grid style={mbgitter,line width=0.3pt},major grid style={mbgitter},xlabel={$x$},ylabel={$y$},x label style={anchor=west},y label style={anchor=south}]\addplot[black,domain=-2:5,samples=120] {((x-2)^2)-2};\node[font=\small,fill=white,inner sep=1pt,anchor=south west] at (axis cs:4.44,3.9536) {$p$};\end{axis}\end{tikzpicture}\par
\end{pfnr}
\begin{pfnr}{}{*a)}{2 BE}
Die lineare Funktion $f$ hat die Gleichung $f(x) = -2x + 2$. Zeichne ihren Graphen in das Koordinatensystem.
\rechenplatz{3}
\end{pfnr}
\begin{pfnr}{}{b)}{2 BE}
Prüfe mit einer Rechnung, ob der Punkt $P(-4\,|\,10)$ auf dem Graphen von $f(x) = -2x + 2$ liegt.
\rechenplatz{3}
\end{pfnr}
\begin{pfnr}{}{c)}{1 BE}
Gib den Scheitelpunkt von $p$ an: $S(\ \ |\ \ )$.
\rechenplatz{3}
\end{pfnr}
\begin{pfnr}{}{*d)}{3 BE}
Eine zweite Parabel $q$ hat die Gleichung $q(x) = -(x-2)^2 - 2$. Gib an, wie viele Punkte $p$ und $q$ gemeinsam haben. Begründe ohne zu rechnen.
\rechenplatz{3}
\end{pfnr}
\end{document}